\section{Conclusione}\label{sec:conclusion}

Questo rapporto ha indagato l'affidabilità delle statistiche ufficiali sull'emigrazione italiana convalidando i dati del registro AIRE rispetto ai record di immigrazione compilati da otto paesi di destinazione europei.
La nostra analisi ha affrontato due domande fondamentali: se i volumi assoluti degli emigranti coincidano e se la distribuzione temporale dei flussi corrisponda tra le due fonti.
L'evidenza statistica è inequivocabile.
Il test dei ranghi con segno di Wilcoxon ha confermato una grave e sistematica sottostima nei dati italiani, dimostrando che l'ISTAT cattura costantemente solo una frazione del vero flusso migratorio.
Inoltre, il test del Chi Quadrato ha dimostrato che i dati ISTAT non riescono a tracciare la corretta tendenza temporale di queste partenze, poiché i ritardi burocratici scollegano le registrazioni AIRE dal momento effettivo della migrazione.
La dipendenza strutturale da atti amministrativi volontari e autodichiarati rende il registro AIRE fondamentalmente incompleto.
Per comprendere veramente la portata, la traiettoria e l'impatto dell'esodo italiano in corso, la contabilità demografica deve guardare oltre i propri confini.
Adottando le statistiche speculari, possiamo iniziare a quantificare le popolazioni invisibili che i registri tradizionali lasciano indietro.
